\documentclass[IB]{TFUOC}
\usepackage[tabla]{tfm-estilo}
\usepackage{tabularx}
\title{T}\titcrt{T}\author{A}\date{2026}\nomPDC{X}\nomPRA{Y}\titulac{M}\area{A}\idioma{Castellano}\credits{15}\parcla{a}\licenc{ccBy}
\abstractidioma{a}\abstractenglish{b}
\begin{document}
\chapter{Árboles de decisión y XGBoost}
Ejemplo de los componentes \emph{ArbolDecision} y \emph{XGBoost} (datos de ejemplo). Los árboles pueden generarse con \texttt{python/arboles.py}.
\begin{figure}[htbp]\centering
\begin{forest} tfmarbol
[{anchura libre $\leq$ 1,80\,m}\arbolinfo{n = 4\,812}, camino
  [{pendiente long. \textgreater\ 8\,\%}\arbolinfo{n = 1\,930}, camino
    [{No accesible}\arbolinfo{n = 640}, hojaclase=no-accesible, camino]
    [{Deficiente}\arbolinfo{n = 1\,290}, hojaclase=deficiente]]
  [{obstáculos $\leq$ 1}\arbolinfo{n = 2\,882}, ausente
    [{Accesible}\arbolinfo{n = 2\,104}, hojaclase=accesible]
    [{Practicable}\arbolinfo{n = 778}, hojaclase=practicable]]]
\end{forest}
\caption{Árbol de decisión para clasificar la accesibilidad de un tramo; se resalta el camino de una observación (datos de ejemplo).}\label{fig:arbol}
\end{figure}

\begin{figure}[htbp]\centering
\etiquetasramas{sí}{no}
\begin{forest} tfmarbol
[{f12 \textless\ 0,35}
  [{f3 \textless\ 1,2}, ausente [{+0,214}] [{$-$0,087}]]
  [{$-$0,152}]]
\end{forest}
\caption{Primer árbol del ensamblado XGBoost; la rama marcada «ausente» recibe los valores nulos (datos de ejemplo).}\label{fig:arbol-xgb}
\end{figure}

\begin{figure}[htbp]\centering
\esquemaboosting[arboles=4, eta={\eta = 0{,}1}]
\caption{Esquema de boosting: cada árbol aprende el residuo de los anteriores y las salidas se suman ponderadas por $\eta$.}\label{fig:boosting}
\end{figure}

\begin{table}[htbp]\centering\small
\caption{Configuración del modelo XGBoost (datos de ejemplo).}\label{tab:config-xgb}
\begin{tabularx}{\linewidth}{@{}l>{\raggedright\arraybackslash}Xll@{}}
\toprule
\cabeceratabla \textbf{Hiperparámetro} & \textbf{Descripción} & \textbf{Rango explorado} & \textbf{Valor} \\
\midrule
\filagrupo{4}{Estructura del árbol}
\hiperparam{max_depth} & Profundidad máxima & 3--10 & \mejor{6} \\
\hiperparam{min_child_weight} & Peso mínimo por hoja & 1--10 & 3 \\
\filagrupo{4}{Aprendizaje}
\hiperparam{learning_rate} ($\eta$) & Tasa de aprendizaje & 0,01--0,3 & 0,05 \\
\bottomrule
\end{tabularx}
\end{table}

\begin{figure}[htbp]\centering
\begin{tikzpicture}
\begin{axis}[tfmimportancia, width=0.8\linewidth, height=5cm, symbolic y coords={Anchura,Pendiente,Obstáculos,Pavimento}]
\addplot coordinates {(0.42,Anchura) (0.27,Pendiente) (0.18,Obstáculos) (0.09,Pavimento)};
\end{axis}
\end{tikzpicture}
\begin{tikzpicture}
\begin{axis}[tfmaprendizaje, width=0.8\linewidth, height=5cm, ylabel={log-loss}]
\addplot+[mark=none] coordinates {(0,1.1) (100,0.6) (200,0.45) (300,0.40) (400,0.41) (500,0.43)};
\addplot+[mark=none] coordinates {(0,1.1) (100,0.5) (200,0.30) (300,0.20) (400,0.14) (500,0.10)};
\legend{Validación, Entrenamiento}
\paradatemprana{312}{parada temprana\\(mejor iteración: 312)}
\end{axis}
\end{tikzpicture}
\caption{Importancia de variables por ganancia y curva de aprendizaje con parada temprana (datos de ejemplo).}\label{fig:xgb-resultados}
\end{figure}
\end{document}
